\chapter{Iteration sequences associated to even integer polynomials}

% The counter is set so that the two lemmas below carry the numbers 2.1, 2.2 of the paper.
\setcounter{theorem}{-1}

In this section, let $g \in \ZZ[X^2]$ be an even polynomial with integer coefficients. Put
$\gamma_1 := \pm g(0)$, and for $n \ge 1$, $\gamma_{n+1} := g(\gamma_n)$. Assume that all
$\gamma_n \ne 0$ and set $\beta_n := \prod_{d \dvd n} \gamma_d^{\mu(n/d)}$ for $n \ge 1$.

\begin{definition}[The sequences $\gamma$ and $\beta$]
  \label{def:gammaseq}
  \lean{QuadraticIterates.EvenPoly, QuadraticIterates.evenPoly_iff_comp_neg_X,
    QuadraticIterates.EvenPoly.eval_neg, QuadraticIterates.EvenPoly.eval_congr,
    QuadraticIterates.EvenPoly.dvd_eval_sub, QuadraticIterates.EvenPoly.map,
    QuadraticIterates.EvenPoly.iterate_eval_neg,
    QuadraticIterates.gammaSeq, QuadraticIterates.gammaSeq_zero,
    QuadraticIterates.gammaSeq_one, QuadraticIterates.gammaSeq_succ,
    QuadraticIterates.gammaSeq_add, QuadraticIterates.map_gammaSeq,
    QuadraticIterates.intCast_gammaSeq, QuadraticIterates.betaSeq,
    QuadraticIterates.betaSeq_eq_moebiusFactorR, QuadraticIterates.intCast_betaSeq}
  \leanok
  $\gamma_1 = \varepsilon\,g(0)$ with $\varepsilon = \pm 1$, $\gamma_{n+1} = g(\gamma_n)$ for
  $n \ge 1$, and $\beta_n = \prod_{d \dvd n} \gamma_d^{\mu(n/d)}$.
  \formalnote{Being even is defined as membership in $R[X^2]$, i.e.\ $g = h(X^2)$ for some $h$;
  over a domain of characteristic $\ne 2$ this is equivalent to $g(-X) = g$
  (\lname{evenPoly\_iff\_comp\_neg\_X}), and both forms are used. As with $c$, the sequence
  $\gamma$ is total with junk value $\gamma_0 = 0$, and $\beta_n$ is an element of the
  coefficient ring, defined as the preimage of the Möbius product in the fraction field. The
  sequences are set up over an arbitrary commutative semiring, and each result below is stated
  over the weakest structure that carries it: a commutative ring for the congruences, a GCD
  domain for strong divisibility, a UFD for the valuation shape, and $\ZZ$ for Lemmas~2.1
  and~2.2. Since the whole construction commutes with ring homomorphisms
  (\lname{map\_gammaSeq}), congruences mod $m$ are proved by computing in $\ZZ/m\ZZ$, and $c$ is
  literally the case $g = X^2 + a$, $\varepsilon = -1$.}
\end{definition}

\begin{lemma}
  \label{lem:2.1}
  \uses{def:gammaseq}
  \lean{QuadraticIterates.not_isSquare_betaSeq, QuadraticIterates.gammaSeq_mul_eq_two_mul,
    QuadraticIterates.prod_gammaSeq_mul_eq, QuadraticIterates.prod_gammaSeq_mul_eq_neg_prod,
    QuadraticIterates.isUnit_prod_gammaSeq_mul, QuadraticIterates.gammaSeq_gcd,
    QuadraticIterates.gammaSeq_associated_gcd, QuadraticIterates.gammaSeq_associated_one,
    QuadraticIterates.gammaSeq_eq_gammaSeq_one}
  \leanok
  Suppose that for all $n \ge 1$ there is an $m$ such that $m \dvd \gamma_n + \gamma_{2n}$, $m$
  is prime to $\gamma_n$, and $-1$ is not a square mod $m$. Then for all $n \ge 2$, $\beta_n$ is
  not a square in $\QQ$.
\end{lemma}
\begin{proof}
  \leanok
  \uses{def:bseq, lem:moebiussign, lem:strongdvd, lem:zmodsquare}
  Let $n > 1$ have prime decomposition $n = p_1^{e_1} \cdots p_r^{e_r}$, where $r \ge 1$, the
  $p_j$ are distinct primes, and all $e_j \ge 1$. Set $n' := p_1 \cdots p_r > 1$ and
  $k := n/n'$. Put $\alpha := \gamma_{2k}$. Let $m$ be such that
  $m \dvd \gamma_k + \gamma_{2k}$, $m$ is prime to $\gamma_k$ and hence to $\alpha$, and $-1$ is
  not a square mod $m$. Then $\gamma_k \equiv -\alpha \bmod m$, and
  $\gamma_{3k} = g_k(\gamma_{2k}) = g_k(-\gamma_k) = g_k(\gamma_k) = \gamma_{2k} = \alpha
  \bmod m$ --- writing $g_k$ for the $k$-fold iterate of $z \mapsto g(z)$, which is an even
  function --- and by induction $\gamma_{tk} \equiv \alpha \bmod m$ for all $t \ge 2$. Therefore
  \[ \beta_n = \prod_{d \dvd n} \gamma_d^{\mu(n/d)}
             = \prod_{t \dvd n'} \gamma_{kt}^{\mu(n'/t)}
             \equiv (-1)^{\mu(n')} \prod_{t \dvd n'} \alpha^{\mu(n'/t)}
             = -1 \bmod m \]
  (for $\mu(n') = \pm 1$ and $\sum_{t \dvd n'} \mu(n'/t) = 0$). Since $-1$ is not a square mod
  $m$, $\beta_n$ cannot be a square in $\QQ$.
  \formalnote{$n'$ is $\rad(n)$, and the Möbius product is split along the partition of the
  divisors of $n'$ into those with $\mu = +1$ and those with $\mu = -1$
  (Lemma~\ref{lem:moebiussign}), which have equal cardinality; $\beta_n$ is then the quotient
  $P/Q$ of the two subproducts, and the computation above becomes $P \equiv -Q \bmod m$ with $Q$
  a unit mod $m$. That a quotient of two such integers cannot be a square when $-1$ is not a
  square mod $m$ is \lname{ZMod.isSquare\_neg\_one\_of\_isSquare\_div}
  (Lemma~\ref{lem:zmodsquare}). The sign $(-1)^{\mu(n')}$ comes from the single index $t = 1$,
  where $\gamma_k \equiv -\alpha$ rather than $\alpha$; this bookkeeping is
  \lname{prod\_gammaSeq\_mul\_eq}.}
\end{proof}

We will now give a handier criterion for the assumptions of Lemma~\ref{lem:2.1} to hold.

\begin{lemma}
  \label{lem:2.2}
  \uses{def:gammaseq, lem:2.1}
  \lean{QuadraticIterates.not_isSquare_betaSeq_of_pos,
    QuadraticIterates.not_isSquare_betaSeq_of_pos_of_not_isSquare_neg_one,
    QuadraticIterates.not_isSquare_betaSeq_of_pos_of_eval_one_emod_four_eq_two,
    QuadraticIterates.not_isSquare_betaSeq_of_pos_of_eval_one_emod_four_eq_three,
    QuadraticIterates.gammaSeq_one_eq_one_of_pos,
    QuadraticIterates.gammaSeq_add_succ_dvd, QuadraticIterates.isCoprime_gammaSeq_add_succ,
    QuadraticIterates.gammaSeq_add_succ_dvd_sub,
    QuadraticIterates.gammaSeq_add_succ_zmod_four_eq_three,
    QuadraticIterates.gammaSeq_add_succ_zmod_eight_eq_six,
    QuadraticIterates.gammaSeq_eq_ite_even, QuadraticIterates.gammaSeq_add_succ_eq_three,
    QuadraticIterates.gammaSeq_eq_eval_one,
    QuadraticIterates.gammaSeq_add_succ_eq_two_mul_eval_one}
  \leanok
  If all the $\gamma_n$ are positive, and
  \begin{enumerate}
    \item[a)] $g(0) = 1$ and $g(1) \equiv 2 \bmod 4$, or
    \item[b)] $g(0) = \pm 1$ and $g(1) \equiv 3 \bmod 4$,
  \end{enumerate}
  then $\beta_n$ is not a square in $\QQ$ for $n \ge 2$.
\end{lemma}
\begin{proof}
  \leanok
  \uses{lem:zmodsquare}
  Notice first that since $\gamma_n + \gamma_{n+1} \dvd \gamma_n + \gamma_{2n}$ (same argument as
  in the proof of Lemma~\ref{lem:2.1}, applied $(n-1)$ times), and $\gamma_n$ and
  $\gamma_{n+1}$ are coprime (since $\gamma_{n+1} = g(\gamma_n) \equiv g(0) = \pm 1
  \bmod \gamma_n$), we can take $m := \gamma_n + \gamma_{n+1}$ in Lemma~\ref{lem:2.1}, if $-1$ is
  not a square mod $m$. This holds for $m \equiv 3 \bmod 4$ and for $m \equiv 6 \bmod 8$.

  a) $\gamma_n \equiv 1$ or $2 \bmod 4$, according to whether $n$ is odd or even, so for all $n$,
  $\gamma_n + \gamma_{n+1} \equiv 3 \bmod 4$.

  b) $\gamma_1 + \gamma_2 \equiv 0 \bmod 4$, and for $n \ge 2$,
  $\gamma_n + \gamma_{n+1} \equiv 6 \bmod 8$, since all $\gamma_n$ $(n \ge 2)$ are
  $\equiv 3 \bmod 4$ and equal mod $8$.
  \formalnote{Positivity of the $\gamma_n$ is what makes $m := \gamma_n + \gamma_{n+1}$ a
  \emph{positive} integer, so that ``$-1$ is not a square mod $m$'' is a statement about
  $\ZZ/m\ZZ$. The two congruence computations are done in $\ZZ/4\ZZ$ resp.\ $\ZZ/8\ZZ$, over an
  arbitrary ring: for a) the only input is $4 = 0$, $g(0) = 1$, $g(1) = 2$, which forces the
  sequence to alternate between $1$ and $2$; for b) the only input is that the relevant residues
  square to $1$, which forces $\gamma_n = g(1)$ for all $n \ge 2$. The case $n = 1$ of b) is the
  separate observation that $4 \dvd \gamma_1 + \gamma_2$ and that $-1$ is not a square mod $4$.}
\end{proof}

In the next section, we will apply this result to the Galois groups of Chapter~1.
